
\chapter{Online Decision Algorithms}

\section*{What is an Online Algorithm?}

An online algorithm is an algorithm that processes its input piece-by-piece in a serial fashion,
 without having the entire input available from the start. In some problems, this doesn't matter -- 
 the obvious solution is online since it requires a single pass over the array only. But in some problems, 
 the online requirement makes the problem much harder and forces us to make decisions without knowing the future.

Online decision algorithms have randomness involved in their analysis for non-adversarial\footnote{They 
are often called oblivious adversaries} inputs. 
Further, the goal of the algorithm usually is not deterministic optimality, but rather to maximise the probability of 
success or to achieve a good competitive ratio\footnote{Competitive ratio is a measure of how well an online algorithm 
performs compared to an optimal offline algorithm that has access to the entire input from the start. The term n-competitive 
is used when an algorithm achieves at least $\frac{1}{n}$ times the "performance" of an optimal offline algorithm.}.

\section{The Dating Problem}

\textbf{The Problem}
Consider the problem
of finding the best partner. Assume that the number of people you will eventually date is fixed, i.e. some known value \(n\). The people appear in a uniformly random order. After dating a person, you must immediately decide whether to commit to them forever or reject them permanently and continue searching. Figure out a strategy that maximises the probability of ending up with the best partner.\footnote{This is a different problem from maximising the expected ``value'' of your partner. Here the objective is specifically to maximise the probability of selecting the single best partner among all \(n\) people.}

We assume that while dating person \(i\), we only know how they compare relative to the people we have already dated. 
So we could sort the partners till now in a total order if needed.\\


\noindent\textbf{Solution:}
\noindent Most course materials, internet reels or shorts I have encountered this problem in 
directly jump to the strategy of:
\begin{quote}
``Reject the first \(r\) people, then accept the first person afterwards who is better than everyone seen so far.''
\end{quote}

\noindent However, they often omit the proof of why an optimal strategy must necessarily have this threshold form. We first prove this fact.

\subsection*{Optimality Proof}

Suppose we are currently considering the \(k\)-th person.

If the current person is \emph{not} the best among the first \(k\) people seen so far, then accepting them can never succeed:
This is because there already exists an earlier person who is strictly better, and that earlier person has already been rejected.
Hence an optimal strategy should only ever consider accepting people who are the best seen so far.\\


Now suppose the current person \emph{is} the best among the first \(k\) people seen so far. 
Let $\bm{q_k}$ be the optimal probability of eventually selecting the globally best person if we reject the $k$-th person and continue optimally afterwards.\\


\textbf{Q] If we instead accept the current person immediately, what is the probability of success?}\\


A] Since the ordering of the \(n\) people is uniformly random, each of the first \(k\) positions is equally likely to contain the globally best person. 
Conditioned on the fact that the current person is the best among the first \(k\), the probability that they are actually 
the globally best person equals $\frac{k}{n}$.\\

\vspace{0.4em}
\noindent Therefore, accepting gives success probability \(\frac{k}{n}\), rejecting gives success probability \(q_k\).
The optimal decision is to accept iff $\frac{k}{n} \ge q_k$\footnote{This is a 
common theme in online algorithms - the optimal decision at each step 
is determined by comparing the expected "value" of actions and choosing 
the one with maximum value. Here our "value" is the success probability. In 
game theory this is called "utility" instead.}. \\


\noindent Now observe that \(\frac{k}{n}\) is \textbf{increasing} with \(k\). \\

\noindent On the other hand, \(q_k\) is \textbf{decreasing} with \(k\). To see this, notice that from time \(k\), one valid strategy is:
\begin{quote}
``Reject one more person no matter what, and then behave optimally from time \(k+1\) onwards.''
\end{quote}

This strategy achieves success probability exactly \(q_{k+1}\). Since \(q_k\) is the \emph{optimal} achievable success probability starting from time \(k\), we must have:
\[
\mathbf{q_k \ge q_{k+1}}
\]

Since \(\frac{k}{n}\) is increasing and \(q_k\) is decreasing, the optimal decision changes at most once from rejecting to accepting as \(k\) increases.
Therefore there exists some threshold \(0 \le r \le n\) such that:
\begin{itemize}
    \item for \(k \le r\), rejecting is optimal,
    \item for \(k > r\), accepting the first ``best-so-far'' person is optimal.
\end{itemize}

\noindent Hence every optimal strategy must be a threshold strategy.
Note that $r=0$ corresponds to accepting the first person, and $r=n$ corresponds to rejecting everyone and ending up with no partner. 
Both of these are valid threshold strategies, though they are not very good ones\footnote{Though it may be argued that having no partner is the best choice - 
this is beyond the scope of this text and is left as an exercise for the reader.}.

\subsection*{Finding the Optimal Threshold}

We now analyse the threshold strategy:
\begin{quote}
``Reject the first \(r\) people, then accept the first person afterwards who is better than everyone seen so far.''
\end{quote}

\noindent Suppose the globally best person appears at position \(k\). For our strategy to succeed:
\begin{enumerate}
    \item We must have \(k > r\)
    \item Among the first \(k-1\) people, the best person must appear within the first \(r\) positions.
\end{enumerate}

\noindent The first condition ensures that we do not automatically reject the best person.
The second condition ensures that no earlier candidate after position \(r\) triggers acceptance before the globally best person appears.

\vspace{0.4em}
\[
\Pr[\text{best person appears at position }k] = \frac{1}{n}
\]
\vspace{0.2em}
Further, among the first \(k-1\) people, each position is equally likely to contain their maximum. 
\vspace{0.2em}
\[
\implies
\Pr[\text{maximum among first }k-1\text{ lies in first }r]
=
\frac{r}{k-1}
\]

\[
\implies
\Pr[\text{success}]
=
\sum_{k=r+1}^{n}
\frac{1}{n}\cdot\frac{r}{k-1}
=
\frac{r}{n}
\sum_{k=r}^{n-1}\frac{1}{k}
=
\frac{r}{n} \left(H(n-1)-H(r-1)\right)
\]
\vspace{0.2em}
\noindent Using the approximation $H(x) \approx \ln(x)$,
\[
\frac{r}{n} \left(H(n-1)-H(r-1)\right)
\approx
\frac{r}{n} \ln\left(\frac{n-1}{r-1}\right) 
\approx
\frac{r}{n} \ln\left(\frac{n}{r}\right)
\]

\[
x=\frac{r}{n} 
\implies
P(x)=x\ln\left(\frac{1}{x}\right)
\]

\[
\implies
P'(x)=\ln\left(\frac{1}{x}\right)-1
\]

\[
P'(x_o)=0
\implies
\ln\left(\frac{1}{x_o}\right)=1
\implies
x_o=\frac{1}{e}
\]
\vspace{0.4em}
\noindent Since $\frac{r_o}{n}=\frac{1}{e} \implies r_o=\frac{n}{e}$, the optimal strategy is --
\begin{quote}
``Reject approximately the first \(\frac{n}{e}\) people, then accept the first person 
afterwards who is better than everyone seen so far.''
\end{quote}

\noindent The corresponding success probability is $\frac{1}{e}\ln e = \frac{1}{e}$ which is 
approximately \(\mathbf{36.8\%}\). Not bad at all!\\


\noindent Strictly speaking, we should also verify that this critical point indeed corresponds to a 
maximum rather than a minimum. One way is via the second derivative test --
\[
P''(x)=-\frac{1}{x}<0\text{,} \hspace{0.5em} \forall \hspace{0.3em} x \in \mathbb R^{>0}
\]
so \(P(x)\) is concave and hence the critical point is necessarily a global maximum.\\


\noindent Alternatively, we can avoid calculus entirely by inspecting the behaviour near the boundaries. 
\[
\because P(x)=x\ln\left(\frac1x\right)
\]

As $x \to 0,P(x) \to 0$ since \(x\) decays faster than \(\ln(1/x)\) grows. \\


\noindent Similarly, as $x \to 1,P(x) \to 0$ because $\ln(1)=0$.\\


\noindent Since \(P(x)\) is positive for \(0<x<1\), rises away from \(0\), and eventually 
falls back to \(0\), the unique critical point at
\[
x=\frac1e
\]
must correspond to the global maximum.

\subsection*{Behaviour for Small \(n\)}

We used some limiting n approximations in finding the optimal threshold (I hope you noticed!).
In case you were not planning on dating a tending-to-infinite number of people, you might be interested in the optimal strategy for small values of \(n\).
For finite \(n\), the exact success probability for threshold \(r\) is:
\[
\max_{0\le r\le n} \left(\frac{r}{n}\sum_{k=r}^{n-1}\frac1k\right)
\]

Computationally we can now easily find max for each \(n\):
\[
\begin{array}{c|cccccc}
n & 2 & 3 & 4 & 5 & 6 & 7\\
\hline
n/e
& 0.74
& 1.10
& 1.47
& 1.84
& 2.21
& 2.58
\\
\hline
\operatorname{round}(n/e)
& 1
& 1
& 1
& 2
& 2
& 3
\\
\hline
\text{Actual optimal }r
& \mathbf{1}
& \mathbf{1}
& \mathbf{1}
& \mathbf{2}
& \mathbf{2}
& \mathbf{3}
\end{array}
\]

\vspace{1em}

The corresponding optimal success probabilities are:

\[
\begin{array}{c|cccccc}
n & 2 & 3 & 4 & 5 & 6 & 7\\
\hline
\Pr[\text{success}]
&
\mathbf{0.500}
&
\mathbf{0.500}
&
\mathbf{0.458}
&
\mathbf{0.433}
&
\mathbf{0.428}
&
\mathbf{0.414}
\end{array}
\]

We observe that the optimal threshold follows the value suggested by
$\frac{n}{e}$ even for relatively small values of \(n\), and it will get closer and closer as \(n\) increases. 
The success probability is already around \(41.4\%\) for \(n=7\), and it will approach the limiting value of \(\frac{1}{e}\approx 36.8\%\) as \(n\) grows.

\vspace{1em}
\section{The Parallel Processing Problem}

Consider a setting with $m$ identical machines processing incoming tasks. The tasks arrive 
one after another, and the arrival times are sufficiently close that 
we may assume they arrive within negligible time of each other.

Each incoming task must immediately be assigned to one of the $m$ machines. Once a decision is made, 
it cannot be reversed. Thus, the algorithm has no knowledge of future tasks while making its decision.\\

\noindent The tasks assigned to a machine are processed sequentially and hence form a queue at that machine.
Suppose the incoming tasks have processing times $T_1,T_2,\dots,T_n$ where $n$ itself is not known in advance.\\


\noindent The objective is to minimize the \textit{makespan}, i.e. the finishing time of the busiest machine:
\[
C_{\max}=\max_{1\le i\le m} T_i
\]
where $T_i$ denotes the total processing time assigned to machine $i$. 
This is a natural objective since the overall processing time
 is determined by the machine that finishes last\footnote{This 
 concept is similar to the Rate Determining Step (RDS) in the field of Chemical Kinetics}.

\subsection{Greedy Algorithm}

A natural strategy is the following greedy algorithm.
\begin{center}
\begin{minipage}{0.9\linewidth}
\begin{framed}
\begin{Verbatim}
Assign(Task T)
{
    find machine with min load
    assign T to this machine
}
\end{Verbatim}
\vspace{-4mm}
\end{framed}
\end{minipage}
\end{center}

\subsubsection*{Performance Analysis}

Interestingly, we do not actually have a polynomial time algorithm to find the optimal \emph{offline} solution to this problem.\\


\noindent We will therefore use lower bounds on $OPT$ (the optimal offline makespan) to derive an upper 
bound on the competitive ratio of the greedy algorithm.\\

\noindent Let $ALG$ denote the makespan produced by the greedy algorithm. Consider the task that 
finishes last in the greedy schedule. This may or may not be the last \emph{assigned} task. Let its 
processing time be $\mathbf{p}$.\\

\noindent Immediately before assigning this job, suppose the minimum machine load was $L$.
By our hypothesis that this machine will be the one that finishes last, the greedy makespan is 
\[
ALG=L+p
\]
\vspace{0.5em}
Hence the total load already assigned before this final job was at least $\mathbf{mL}$.
Including the final job, the total processing load is therefore at least $\mathbf{mL+p}$.\\


\noindent By Pigeonhole Principle, at least one machine must have load at least the average load, so --
\[
OPT\ge \frac{mL+p}{m} = L+\frac{p}{m}
\]

Since any valid schedule must process the job of size $p$,
\[
OPT\ge p \implies OPT(1-\frac{1}{m})\ge p\left(1-\frac{1}{m}\right)
\]

Adding the two inequalities gives:
\[
OPT(2-\frac{1}{m})\ge L+p = ALG
\implies
\boxed{\frac{ALG}{OPT}\le 2-\frac1m}
\]
\vspace{0.5em}
Therefore the greedy algorithm is $\left(2-\frac1m\right)$-competitive, sometimes written as $2$-competitive 
since the difference is negligible for large $m$.

\subsection{Tightness of the Bound}

The previous analysis showed that the greedy algorithm is $\left(2-\frac1m\right)$-competitive.
 We now construct an example showing that the analysis cannot,
 in general, be improved beyond a factor of $\left(2-\frac1m\right)$. \\

Consider $m=4$ machines and the following sequence of incoming jobs:
\[
1,1,1,\dots,1 \text{ (12 times)},4
\]

The greedy algorithm will assign the first 12 jobs as follows:
\[
\begin{array}{c|cccc}
\text{Machine }i & 1 & 2 & 3 & 4\\
\hline
\text{Jobs assigned} & 1,1,1 & 1,1,1 & 1,1,1 & 1,1,1
\end{array}
\]

After processing these 12 jobs, each machine has load 3. The last job of size 4 will be assigned to any machine (say machine 1), resulting in a final makespan of 7.
On the other hand, an optimal offline algorithm could have assigned the first 12 jobs as follows:
\[
\begin{array}{c|cccc}
\text{Machine }i & 1 & 2 & 3 & 4\\
\hline
\text{Jobs assigned} & 1,1,1,1 & 1,1,1,1 & 1,1,1,1 & 4
\end{array}
\]

After processing these 13 jobs, each machine has load 4. Thus the optimal makespan is 4,
 and the competitive ratio is $\frac{7}{4}$, which is indeed our bound of $2-\frac{1}{4}$.

